% ---------------------------------------------------------------
% CHAPTER 4: PROJECT SIMULATION: SKYFOX CINEMAS
% ---------------------------------------------------------------
\chapter{Project Simulation: SkyFox Cinemas}

\section{Project Overview}

\textit{SkyFox Cinemas} constitutes a comprehensive internal operational and 
reservation management system engineered explicitly for cinema operators. Contrasting 
with third-party marketplace mediators, SkyFox delivers enterprise resource management 
capabilities, furnishing proprietors with unrestricted command of resources, operational 
schedules, and personnel administration. The system was constructed through \textbf{three 
sequential Agile delivery iterations} (seven-day cycles) executed by a multidisciplinary 
cohort of eleven interns. The author facilitated technology oversight (\textbf{Tech Lead (TL)}) 
throughout all iterations.

\begin{figure}[H]
\centering
\begin{tikzpicture}[node distance=0.45cm]
  \node[rectangle, fill=idfcblue, text=white, minimum width=12.5cm,
        minimum height=0.8cm, font=\bfseries\large, rounded corners=4pt]
    (title) {SkyFox Cinemas --- Project Blueprint};
  \node[block, below=0.4cm of title, xshift=-3.2cm, text width=5.5cm] (p1)
    {Platform: Internal Cinema ERP \& Booking System};
  \node[block, below=0.4cm of title, xshift= 3.2cm, text width=5.5cm] (p2)
    {Team Size: 11 Interns};
  \node[block, below=0.4cm of p1, text width=5.5cm] (p3)
    {Methodology: Agile/Scrum (3 $\times$ 1-week sprints)};
  \node[block, below=0.4cm of p2, text width=5.5cm] (p4)
    {Author's Role: Tech Lead (all 3 sprints)};
  \node[block, below=0.4cm of p3, text width=5.5cm] (p5)
    {Backend: GoLang + Gin Framework + Goose Migrations};
  \node[block, below=0.4cm of p4, text width=5.5cm] (p6)
    {Frontend: React JS};
  \node[block, below=0.4cm of p5, text width=5.5cm] (p7)
    {Database: PostgreSQL};
  \node[block, below=0.4cm of p6, text width=5.5cm] (p8)
    {DevOps: Bitbucket + CI/CD Pipelines};
\end{tikzpicture}
\caption{SkyFox Cinemas --- Project Details}
\end{figure}

\section{Sprint Breakdown}

\begin{figure}[H]
\centering
\begin{tikzpicture}[node distance=0.5cm and 0.8cm]
  % Sprint 1
  \node[rectangle, fill=idfcblue!80, text=white, rounded corners=4pt,
        minimum width=3.6cm, minimum height=0.7cm, font=\bfseries\small]
    (s1h) {SPRINT 1 --- Week 1};
  \node[block, below=0.3cm of s1h, minimum width=3.6cm,
        text width=3.4cm, font=\footnotesize, minimum height=1.8cm] (s1b)
    {Project setup, DB schema design, movie listing API, React scaffolding};
  % Sprint 2
  \node[rectangle, fill=idfcblue!80, text=white, rounded corners=4pt,
        minimum width=3.6cm, minimum height=0.7cm, font=\bfseries\small,
        right=0.85cm of s1h] (s2h) {SPRINT 2 --- Week 2};
  \node[block, below=0.3cm of s2h, minimum width=3.6cm,
        text width=3.4cm, font=\footnotesize, minimum height=1.8cm] (s2b)
    {Seat selection UI, booking API with concurrency handling, RBAC system};
  % Sprint 3
  \node[rectangle, fill=idfcblue!80, text=white, rounded corners=4pt,
        minimum width=3.6cm, minimum height=0.7cm, font=\bfseries\small,
        right=0.85cm of s2h] (s3h) {SPRINT 3 --- Week 3};
  \node[block, below=0.3cm of s3h, minimum width=3.6cm,
        text width=3.4cm, font=\footnotesize, minimum height=1.8cm] (s3b)
    {Partial check-in logic, Admin control APIs, DB migrations via Goose, presentation};
  \draw[arrow] (s1b.east) -- (s2b.west);
  \draw[arrow] (s2b.east) -- (s3b.west);
\end{tikzpicture}
\caption{SkyFox Cinemas --- Agile Sprint Breakdown}
\end{figure}

\section{Technical Contributions}

\subsection{Concurrent Seat Booking System}

The seat reservation subsystem was engineered to accommodate high-volume competing 
attempts where simultaneous customers endeavor to acquire the same seating. The 
implementation employed PostgreSQL transaction isolation (\texttt{SELECT FOR UPDATE}) 
contained within a transactional boundary. This reservation method guaranteed that 
database records were exclusively controlled throughout transaction operation, 
prohibiting simultaneous acquisitions and confirming single-success-only booking 
confirmation.

\subsection{Role-Based Access Control (RBAC) \& Secure Onboarding}
To maintain operational safety, I established a tiered RBAC composition. The arrangement 
encompasses:
\begin{itemize}
    \item \textbf{Authority Distribution:} The Owner manages Admin profile generation with 
    restricted authorization levels, facilitating secure responsibility allocation.
    \item \textbf{Initial Access Protocol:} Novel Admins receive momentary login information. 
    A database indicator (\texttt{is\_first\_login}) activates mandatory security phrase 
    modification, confirming that machine-assigned temporary access tokens are promptly changed.
\end{itemize}

\subsection{Owner-Centric Schedule Management}
A significant architectural differentiator in SkyFox constitutes the operational 
authority distributed to cinema administrators. I engineered an Administration Interface 
facilitating:
\begin{itemize}
    \item \textbf{Movie Assignment Modification:} Altering the content scheduled for 
    particular time slots to enhance monetization per performance metrics.
    \item \textbf{Reschedule/Termination Capabilities:} Transferring existing appointments 
    or terminating presentations with concurrent status alteration, encompassing contention 
    resolution logic to guarantee auditorium accessibility.
\end{itemize}

\subsection{Granular Partial Check-In}
To handle pragmatic group reservation conditions, I constructed a "Granular Check-In" 
capability. This permits a sales assistant to finalize individual attendees from a 
collective reservation (e.g., finalizing 3 guests from a ten-seat allocation) with 
unaddressed participants keeping "Awaiting" classification. This delivers the proprietor 
with granular, present occupancy intelligence.

\subsection{Database Migrations with Goose}
Database structure adjustments, encompassing intricate relationship definitions for 
per-attendee enrollment and \texttt{users} table modifications for authorization levels, 
were overseen via \textbf{Goose}. This approach guaranteed database architecture remained 
version-managed and synchronized among all development systems.

\section{Tech Lead Responsibilities}

\begin{figure}[H]
\centering
\begin{tikzpicture}[node distance=0.45cm]
  \node[rectangle, fill=idfcblue, text=white, minimum width=12.5cm,
        minimum height=0.8cm, font=\bfseries, rounded corners=4pt]
    (title) {Tech Lead Responsibilities Across 3 Sprints};
  \node[block, below=0.4cm of title, xshift=-3.2cm,
        text width=5.5cm, minimum height=1.1cm] (tl1)
    {\textbf{Sprint Planning}\\Translating owner requirements into technical tasks};
  \node[block, below=0.4cm of title, xshift= 3.2cm,
        text width=5.5cm, minimum height=1.1cm] (tl2)
    {\textbf{Security Architecture}\\Designing the RBAC and secure onboarding logic};
  \node[block, below=0.4cm of tl1,
        text width=5.5cm, minimum height=1.1cm] (tl3)
    {\textbf{Code Quality}\\Conducting code reviews to ensure secure credential handling};
  \node[block, below=0.4cm of tl2,
        text width=5.5cm, minimum height=1.1cm] (tl4)
    {\textbf{System Integration}\\Ensuring frontend-backend parity for the Admin dashboard};
\end{tikzpicture}
\caption{Tech Lead Responsibilities in SkyFox Cinemas}
\end{figure}

% ---------------------------------------------------------------
% CHAPTER 5: RESILIENT EVENT CONSUMER FRAMEWORK
% ---------------------------------------------------------------
\chapter{Resilient Event Consumer Framework}

\section{Problem Context}

IDFC FIRST Bank's distributed service architecture deploys Apache Kafka as its principal 
information stream architecture. Multiple microservices transmit categorized information --- 
authentication events, financial operations, customer verification updates --- consumed by 
subsequent processing services, encompassing the Event Record Writer, which documents 
information into PostgreSQL conforming to RBI directives.

\begin{figure}[H]
\centering
\begin{tikzpicture}[node distance=0.6cm]
  \node[medblock] (c) {Consume Event from Kafka};
  \node[medblock, below=of c] (t)
    {Attempt DB Write to PostgreSQL};
  \node[decision, below=of t] (d) {Write\\Successful?};
  \node[greenblock, minimum width=3.8cm,
        below left=0.8cm and 0.5cm of d] (ok)
    {Offset Committed\\$\rightarrow$ Done};
  \node[redblock, minimum width=3.8cm,
        below right=0.8cm and 0.5cm of d] (err)
    {Log Error\\$\rightarrow$ Move On};
  \node[redblock, minimum width=3.8cm, below=0.5cm of err,
        draw=redborder, fill=redbg!60] (lost)
    {\textbf{EVENT LOST}\\(Permanent \& Unrecoverable)};
  \draw[arrow] (c) -- (t);
  \draw[arrow] (t) -- (d);
  \draw[arrow] (d.west) -|
    node[above left, font=\footnotesize\bfseries, text=greenborder]{Yes}
    (ok.north);
  \draw[arrow] (d.east) -|
    node[above right, font=\footnotesize\bfseries, text=redborder]{No}
    (err.north);
  \draw[arrow] (err) -- (lost);
\end{tikzpicture}
\caption{Legacy Consumer --- Critical Reliability Gap}
\end{figure}

\section{Design Goals}

\begin{figure}[H]
\centering
\begin{tikzpicture}[node distance=0.45cm]
  \node[widepurpleblock, text width=12.2cm] (g1)
    {\textbf{Goal 1 --- Zero Event Loss:}
     No information records eliminated during momentary service disruptions};
  \node[widepurpleblock, below=of g1, text width=12.2cm] (g2)
    {\textbf{Goal 2 --- Duplicate Prevention:}
     Duplicate-safe handling conforming to Kafka's minimum-once guarantee};
  \node[widepurpleblock, below=of g2, text width=12.2cm] (g3)
    {\textbf{Goal 3 --- Recoverability:}
     Information surpassing reprocessing maximums is maintained and restartable};
  \node[widepurpleblock, below=of g3, text width=12.2cm] (g4)
    {\textbf{Goal 4 --- Observability:}
     Standardized logs and Prometheus indicators for all fault modes};
  \node[widepurpleblock, below=of g4, text width=12.2cm] (g5)
    {\textbf{Goal 5 --- Reusability:}
     Available as a package applicable to any event-oriented service at the bank};
\end{tikzpicture}
\caption{Framework Design Goals}
\end{figure}

\section{System Architecture}

\begin{figure}[H]
\centering
\begin{tikzpicture}[node distance=0.6cm]
  \node[medblock] (k) {Kafka Event Delivery};
  \node[medblock, below=of k] (i)
    {Idempotency Check (Redis TTL Store)};
  \node[decision, below=of i] (seen) {Already\\Seen?};
  \node[yellowblock, minimum width=4cm,
        below right=0.8cm and 0.3cm of seen] (skip)
    {Skip Silently\\(Duplicate Prevented)};
  \node[medblock, minimum width=4cm,
        below left=0.8cm and 0.3cm of seen] (retry)
    {Retry Engine\\(Exponential Backoff + Jitter)};
  \node[decision, below=0.7cm of retry] (success) {Processed\\OK?};
  \node[greenblock, minimum width=3.8cm,
        below left=0.8cm and 0.2cm of success] (commit)
    {DB Write $\rightarrow$\\Offset Commit};
  \node[redblock, minimum width=3.8cm,
        below right=0.8cm and 0.2cm of success] (dlq)
    {DLQ Publisher\\$\rightarrow$ PagerDuty Alert};
  \node[yellowblock, minimum width=3.8cm, below=0.5cm of dlq] (replay)
    {DLQ Replayer CLI\\(On Demand)};
  \draw[arrow] (k) -- (i);
  \draw[arrow] (i) -- (seen);
  \draw[arrow] (seen.east) -|
    node[above right, font=\footnotesize\bfseries]{Yes} (skip.north);
  \draw[arrow] (seen.west) -|
    node[above left, font=\footnotesize\bfseries]{No} (retry.north);
  \draw[arrow] (retry) -- (success);
  \draw[arrow] (success.west) -|
    node[above left, font=\footnotesize\bfseries, text=greenborder]{Yes}
    (commit.north);
  \draw[arrow] (success.east) -|
    node[above right, font=\footnotesize\bfseries, text=redborder]{Exhausted}
    (dlq.north);
  \draw[arrow] (dlq) -- (replay);
\end{tikzpicture}
\caption{Resilient Event Consumer Framework --- End-to-End Architecture}
\end{figure}

\section{Component 1: Core Retry Engine}

\subsection{Failure Classification}

\begin{figure}[H]
\centering
\begin{tikzpicture}[node distance=0.6cm]
  \node[medblock] (in) {Incoming Processing Failure};
  \node[decision, below=0.5cm of in] (d) {Failure\\Type?};
  \node[redblock, minimum width=4.8cm, text width=4.6cm,
        below left=0.8cm and 1.5cm of d] (trans)
    {\textbf{Transient}\\DB timeout, network blip,\\connection pool exhausted};
  \node[yellowblock, minimum width=4.8cm, text width=4.6cm,
        below right=0.8cm and 1.5cm of d] (perm)
    {\textbf{Permanent}\\Malformed event,\\schema mismatch, invalid payload};
  \node[block, minimum width=4.8cm, below=0.5cm of trans] (ret)
    {Retry: Exponential Backoff\\+ Jitter};
  \node[redblock, minimum width=4.8cm, below=0.5cm of perm] (dlq2)
    {Skip Retries $\rightarrow$ DLQ Directly};
  \draw[arrow] (in) -- (d);
  \draw[arrow] (d.west) -| (trans.north);
  \draw[arrow] (d.east) -| (perm.north);
  \draw[arrow] (trans) -- (ret);
  \draw[arrow] (perm) -- (dlq2);
\end{tikzpicture}
\caption{Failure Classification in the Retry Engine}
\end{figure}

\subsection{Exponential Backoff with Jitter}

Constant-rate reprocessing generates a \textit{coordinated spike} condition at scale. The 
computational formula employed is:
\begin{equation}
  \text{delay}_n = \min\!\left(\text{baseDelay} \times 2^{n-1}
  + \text{random}(0,\, \text{jitter}),\ \text{maxDelay}\right)
\end{equation}
The randomization component desynchronizes reprocessing surges among distributed consumers 
--- the recognized convention implemented by Amazon infrastructure, Google cloud utilities, 
and Netflix fault-tolerance technologies.

\begin{lstlisting}[language=bash, caption={Retry Engine Configuration}]
RETRY_MAX_ATTEMPTS=4
RETRY_BASE_DELAY_MS=2000
RETRY_MAX_DELAY_MS=30000
RETRY_JITTER_MS=500
\end{lstlisting}

\section{Component 2: Redis-Backed Idempotency Store}

Kafka delivers minimum-once transmission assurance. When a consuming process terminates 
post-transmission preservation yet preceding message tracking acknowledgment, renewed 
delivery manifests, producing replicated records. The architecture implements Redis-based 
duplicate elimination.

\begin{figure}[H]
\centering
\begin{tikzpicture}[node distance=0.6cm]
  \node[medblock] (arrive) {Event Arrives from Kafka};
  \node[medblock, below=of arrive] (check)
    {Check Redis: Key = \texttt{audit:processed:\{event\_id\}}};
  \node[decision, below=of check] (d) {Key\\Exists?};
  \node[yellowblock, minimum width=4cm,
        below right=0.8cm and 0.3cm of d] (skip2)
    {Skip Silently\\(Duplicate Prevented)};
  \node[medblock, minimum width=4cm,
        below left=0.8cm and 0.3cm of d] (proc)
    {Process Event\\$\rightarrow$ Write to PostgreSQL};
  \node[greenblock, minimum width=4cm, below=0.5cm of proc] (store)
    {Store Key in Redis\\(TTL: 24 hours)};
  \node[greenblock, minimum width=4cm, below=0.5cm of store] (commit2)
    {Commit Kafka Offset};
  \draw[arrow] (arrive) -- (check);
  \draw[arrow] (check) -- (d);
  \draw[arrow] (d.east) -|
    node[above right, font=\footnotesize\bfseries]{Yes} (skip2.north);
  \draw[arrow] (d.west) -|
    node[above left, font=\footnotesize\bfseries]{No} (proc.north);
  \draw[arrow] (proc) -- (store);
  \draw[arrow] (store) -- (commit2);
\end{tikzpicture}
\caption{Redis-Backed Idempotency Store --- Processing Flow}
\end{figure}

Redis was chosen rather than database queries due to nanosecond-range ($\sim$0.1ms) 
instantaneous lookup response, programmed expiration-based automatic cleanup, and minimal 
supplemental workload upon the operational database.

\section{Component 3: DLQ Publisher and Replayer CLI}

\begin{figure}[H]
\centering
\begin{tikzpicture}[node distance=0.6cm]
  \node[wideredblock] (exhaust)
    {Retry Engine Exhausted --- All 4 Attempts Failed};
  \node[wideblock, below=of exhaust] (enrich)
    {Enrich Event with Metadata:
     failure\_reason, attempt\_count, timestamps, consumer\_id};
  \node[wideblock, below=of enrich] (pub)
    {Publish to Kafka Topic: \texttt{audit.events.dlq}};
  \node[wideyellowblock, below=of pub] (alert)
    {Dispatch PagerDuty + Slack Alert to On-Call Engineer};
  \draw[arrow] (exhaust) -- (enrich);
  \draw[arrow] (enrich) -- (pub);
  \draw[arrow] (pub) -- (alert);
\end{tikzpicture}
\caption{DLQ Publisher --- Enrichment and Alerting Flow}
\end{figure}

A standalone GoLang utility application permits operations staff to replay information from 
the dead-letter repository on-demand once the underlying problem is corrected.

\begin{lstlisting}[language=bash, caption={DLQ Replayer CLI --- Sample Usage}]
$ audit-replayer --topic audit.events.dlq \
                 --from "2026-04-25T10:00:00Z"
[INFO] Found 47 events in DLQ
[INFO] Replaying... 47/47
[OK]   Successfully reprocessed: 45
[WARN] Still failing: 2  (written to replay-failures.json)
\end{lstlisting}

\section{Component 4: Observability Layer}

\begin{figure}[H]
\centering
\begin{tikzpicture}[node distance=0.5cm]
  \node[rectangle, fill=idfcblue, text=white, minimum width=12.5cm,
        minimum height=0.8cm, font=\bfseries, rounded corners=4pt]
    (title) {Observability Layer --- Two Pillars};
  \node[block, below=0.5cm of title, xshift=-3.2cm,
        text width=5.5cm, minimum height=2.5cm] (log)
    {\textbf{Structured JSON Logging}\\[4pt]
     \footnotesize Fields: timestamp, level, event\_id,\\
     attempt, error, retry\_in\_ms, consumer\_id\\[4pt]
     $\rightarrow$ Ingested by Splunk / ELK / Datadog};
  \node[block, below=0.5cm of title, xshift= 3.2cm,
        text width=5.5cm, minimum height=2.5cm] (prom)
    {\textbf{Prometheus \texttt{/metrics} Endpoint}\\[4pt]
     \footnotesize events\_processed (success/failed)\\
     retry\_attempts\_total per level\\
     dlq\_published\_total\\
     p99 processing latency (ms)\\[4pt]
     $\rightarrow$ Grafana Dashboards + Auto-Alerts};
  \node[greenblock, below=0.5cm of log, text width=5.5cm] (g1)
    {Spike in \texttt{dlq\_published\_total}\\triggers PagerDuty before SLA breach};
  \node[greenblock, below=0.5cm of prom, text width=5.5cm] (g2)
    {Machine-parseable; filterable\\by event\_id or consumer\_id};
\end{tikzpicture}
\caption{Observability Layer --- Structured Logging and Prometheus Metrics}
\end{figure}

\section{Component 5: Configuration Management}

\begin{figure}[H]
\centering
\begin{tikzpicture}[node distance=0.6cm]
  \node[wideblock, text width=12.2cm] (env)
    {Load ALL Config from Environment Variables at Startup\\
    (Kafka broker, topics, consumer group, DB host, Redis URL, retry params)};
  \node[wideblock, below=of env, text width=12.2cm] (val)
    {Startup Validation: Ping Kafka Broker, PostgreSQL, Redis};
  \node[decision, below=of val] (ok3) {All Checks\\Pass?};
  \node[greenblock, minimum width=4.2cm,
        below left=0.8cm and 1.2cm of ok3] (start)
    {Service Starts Normally};
  \node[redblock, minimum width=4.2cm,
        below right=0.8cm and 1.2cm of ok3] (abort)
    {Startup ABORTED\\(Prevents Silent Misconfiguration)};
  \draw[arrow] (env) -- (val);
  \draw[arrow] (val) -- (ok3);
  \draw[arrow] (ok3.west) -|
    node[above left, font=\footnotesize\bfseries, text=greenborder]{Yes}
    (start.north);
  \draw[arrow] (ok3.east) -|
    node[above right, font=\footnotesize\bfseries, text=redborder]{No}
    (abort.north);
\end{tikzpicture}
\caption{Configuration Management --- Startup Validation Flow}
\end{figure}

\section{Component 6: Testing Strategy}

\subsection{Unit Tests}

\begin{figure}[H]
\centering
\begin{tikzpicture}[node distance=0.4cm]
  \node[rectangle, fill=idfcblue, text=white, minimum width=12.5cm,
        minimum height=0.7cm, font=\bfseries, rounded corners=4pt]
    (title) {Unit Test Suite --- Key Test Cases};
  \node[widegreenblock, below=0.35cm of title, text width=12.2cm,
        font=\footnotesize] (t1)
    {\texttt{TestRetryEngine\_TransientFailure\_RetriesCorrectly}
     --- Engine retries on transient errors up to max attempts};
  \node[widegreenblock, below=0.3cm of t1, text width=12.2cm,
        font=\footnotesize] (t2)
    {\texttt{TestRetryEngine\_PermanentFailure\_SkipsToDLQ}
     --- Engine bypasses retries for permanent failures};
  \node[widegreenblock, below=0.3cm of t2, text width=12.2cm,
        font=\footnotesize] (t3)
    {\texttt{TestRetryEngine\_ExponentialBackoff\_DelayDoubles}
     --- Delay doubles each attempt within configured bounds};
  \node[widegreenblock, below=0.3cm of t3, text width=12.2cm,
        font=\footnotesize] (t4)
    {\texttt{TestIdempotencyStore\_DuplicateEvent\_Skipped}
     --- Second occurrence of same event\_id is silently skipped};
  \node[widegreenblock, below=0.3cm of t4, text width=12.2cm,
        font=\footnotesize] (t5)
    {\texttt{TestDLQPublisher\_ExhaustedRetries\_PublishesToDLQ}
     --- Event reaches DLQ topic after max retry attempts};
  \node[widegreenblock, below=0.3cm of t5, text width=12.2cm,
        font=\footnotesize] (t6)
    {\texttt{TestConfig\_MissingEnvVar\_StartupFails}
     --- Service refuses to boot when required config is absent};
\end{tikzpicture}
\caption{Unit Test Coverage --- Key Test Cases}
\end{figure}
